\section{当前系统的边界和下一步}

这个项目已经从自己的 ROM 启动，运行 x86-64 C++20 内核，并提供进程、
Thread、虚拟内存、文件系统、Shell、信号、终端与持久化。它仍是一台单 CPU
的教学系统。理解这个边界很重要：现有对象模型为继续扩展提供起点，但并没有
自动获得 Linux 的硬件范围、兼容接口和长期优化。

\topic{与现代 Linux 对齐的是问题拆分方法}
Process 与 Thread 分开，fd 与打开文件实例分开，VMA 与 PTE 分开，VFS 与
具体文件系统分开。这些拆分使共享和释放能够被单独说明。它们不意味着项目
支持 Linux 的全部系统调用、驱动、权限模型或并发规模。

\begin{osgridlongtable}{L{0.22\textwidth}L{0.31\textwidth}L{0.35\textwidth}}
\textbf{后续方向} & \textbf{新增的核心问题} & \textbf{应先补的基础}\\ \hline
SMP & 同一对象被多个 CPU 同时访问、跨核 TLB 失效 & 原子操作、每 CPU
调度队列、IPI 与锁顺序审计\\ \hline
AHCI/NVMe & DMA、描述符队列与并行完成 & PCI 枚举、物理连续缓冲、IOMMU
边界\\ \hline
网络 & 包队列、协议超时、socket 生命周期 & 网卡驱动、定时器扩展、校验和与
缓冲所有权\\ \hline
动态链接 & 多个进程共享代码、重定位和符号解析 & ELF 动态段、共享对象 ABI
与只读页缓存\\ \hline
权限与凭据 & “能找到对象”不再等于“可以操作对象” & uid/gid、mode、
capability 与路径检查时机\\ \hline
swap & 不驻留页面也可能保存着匿名数据 & 页回收、反向映射、I/O 失败与
OOM 策略\\ \hline
\end{osgridlongtable}

\topic{先做 SMP 会改变哪些旧假设}
当前只有一个 BSP，IRQ 仍会打断普通内核代码，但同一时刻不会有两个 CPU
执行普通 Thread。加入第二个 CPU 后，关闭本地中断只能阻止本 CPU 的 IRQ，
不能保护另一个 CPU 正在修改的链表。引用计数、Ready 队列、页表回收和
FileTable 都要重新检查内存顺序与锁粒度。

跨核修改页表还需要 TLB shootdown。释放页帧前，内核必须确认其他 CPU 不再
缓存旧映射；否则旧 TLB 项仍能访问已经交给别人的物理页。这是一项独立工程，
不适合只在启动代码中“打开多个核”就宣布完成。

\topic{换用现代磁盘控制器会引入 DMA}
PIO 由 CPU 逐字读写端口，速度慢，但内存所有权清楚。AHCI 或 NVMe 让设备
根据描述符直接访问 RAM。驱动必须保证描述符和数据缓冲在设备完成前仍存在，
处理物理地址、对齐、缓存一致性、超时取消和迟到中断。设备返回失败时，已经
提交到队列的请求也不能被简单释放。

\topic{一只数据包从网卡到用户缓冲区经过谁}
网卡收到一帧以后，缓冲可能依次属于 DMA ring、驱动、协议栈和 socket 接收
队列。复制可以简化所有权但增加成本；共享可以减少复制却要求引用和回收协议。
第一版网络栈应从固定网卡模型、ARP、IPv4、ICMP 和 UDP 小步推进，每一步都
保留丢包、校验失败与超时测试。

\topic{优化应该由测量触发}
线性 VMA、简单目录查找和全局结构在当前规模下更容易验证。若测量显示路径
解析或页故障查找成为瓶颈，再引入树、哈希或缓存。新的数据结构要说明失效
条件、并发修改和内存压力下的回收，不能只比较一次正常查询的速度。

\topic{选择下一阶段的三个判断}
优先选择能够暴露一个新系统问题、又能用现有观测手段检查的机制。确认它的
最小版本不会同时要求三项尚不存在的基础设施。最后写出明确的退出条件：哪些
成功、边界和故障测试通过以后，这一阶段才停止扩张。这样下一步仍是一段可以
走读的系统演进，而不是功能清单。
